\documentclass[17pt]{article}

\usepackage{cmap}	% поиск по документу
\usepackage{hyperref}	% поддержка гиперссылок
\usepackage{geometry}	% настройка геометрии документа
\usepackage{multicol}   % настройка шаблона расположения текста
\usepackage{extsizes}	% расширенная настройка размеров шрифтов
\usepackage{indentfirst}    % отступ для первого абзаца
\usepackage{tabularx}	% расширенные таблицы
\usepackage{makecell}   % настройка клеток таблиц
\usepackage{caption}    % настройка подписей к вставкам
\usepackage[table]{xcolor}  % покраска таблиц
\usepackage{graphicx}	% вставка изображений
\usepackage{wrapfig, booktabs}	% оформление таблиц
\usepackage{graphicx, caption, subcaption}
\usepackage{minted}    % вставка кода

\usepackage{color}	% настройка цвета текста - \textcolor
\usepackage{soul}	% выделение текста - \hl
\usepackage[T2A]{fontenc}	% поддержка кириллических символов
\usepackage[utf8]{inputenc} % поддержка кириллических символов
\usepackage[english,russian]{babel}	% поддержка кириллических символов
\usepackage{csquotes, biblatex}   % вывод библиографии

\usepackage{mathcmd}	% математические команды
\usepackage{amsfonts}	% математические символы 1
\usepackage{amssymb}	% математические символы 2
\usepackage{amsmath}	% математические символы 3
\usepackage{amsthm}	% работа с теоремами

% Настройка окружения теорем 1 (RU)
\newtheorem{theorem}{Теорема}[section]
\newtheorem{corollary}{Следствие}[theorem]
\newtheorem{lemma}[theorem]{Лемма}
\newtheorem{example}{Пример}[section]
\newtheorem*{remark}{Утверждение}
\addto\captionsrussian{\renewcommand\proofname{Док-во:}}

% Настройка окружения теорем 2 (EN)
%\newtheorem{theorem}{TH}[section]
%\newtheorem{corollary}{CO}[theorem]
%\newtheorem{lemma}[theorem]{LM}
%\newtheorem{example}{EX}[section]
%\addto\captionsrussian{\renewcommand\proofname{PF:}}

\renewcommand{\labelitemii}{$\circ$}
\renewcommand{\labelitemiii}{-}

\addbibresource{biblio.bib}
\graphicspath{ {./images/} }

% Настройка оформления вставок кода
\setminted{xleftmargin=\parindent, fontsize=\footnotesize, linenos, tabsize=2, breaklines}
\usemintedstyle{emacs} 

% Настройка отступов
\geometry{
    a4paper,
    % total={200mm,287mm},
    left=10mm,
    right=10mm,
    top=10mm,
    bottom=20mm
}

% Заголовок
\title{\large{Методы моделирования и анализа стохастических систем\\ Лабораторная работа № 4}}
\author{\large{Костарев Иван, МЕНМ-150201}}
\date{весенний семестр 2025/2026}



\begin{document}

\maketitle
\tableofcontents
\newpage



\section*{Исследование стохастического возбуждения колебаний в модели ФХН}

\section{Сравнение стохастических решений}

Проведем сравнение стохастических решений для $a = 1.1$, $a = 1.01$

(а) Построим графики $x(\varepsilon)$, пропустив переходный процесс $T = 100$ (Рис. \ref{fig:excitement_amplitude}).

\begin{figure}[htbp]
    \centering
    \includegraphics[width=1\textwidth]{excitement_amplitude.png}
    \caption{Амплитуда координаты $x$ при различных значениях $\varepsilon$.}
    \label{fig:excitement_amplitude}
\end{figure}

Локализуем зону стохастического возбуждения на основе 10 симуляций (Рис. \ref{fig:excitement_zone}).

\begin{figure}[htbp]
    \centering
    \includegraphics[width=1\textwidth]{excitement_zone.png}
    \caption{$\varepsilon$-зона стохастического возбуждения.}
    \label{fig:excitement_zone}
\end{figure}


(б) Построим временные ряды при различных значениях $\varepsilon$: с возбуждением и без возбуждения (Рис. \ref{fig:time_series_1_1}, \ref{fig:time_series_1_01}).

\begin{figure}[htbp]
    \centering
    \includegraphics[width=1\textwidth]{time_series_a=1,1.png}
    \caption{Временные ряды для $a=1.1$.}
    \label{fig:time_series_1_1}
\end{figure}

\begin{figure}[htbp]
    \centering
    \includegraphics[width=1\textwidth]{time_series_a=1,01.png}
    \caption{Временные ряды для $a=1.01$.}
    \label{fig:time_series_1_01}
\end{figure}



\newpage
\section{Доверительные эллипсы в допороговой и надпороговой зонах}

Построим доверительные эллипсы в допороговой и надпороговой зонах для $a = 1.1$, $a = 1.01$ (Рис. \ref{fig:confidence_ellipse_1_1}, \ref{fig:confidence_ellipse_1_01}).

Траектория системы демонстрирует ожидаемое поведение и, находясь в режиме возбуждения, выходит за пределы доверительного эллипса.

\begin{figure}[htbp]
    \centering
    \includegraphics[width=1\textwidth]{confidence_ellipse_a=1,1.png}
    \caption{Доверительные эллипсы для $a=1.1$.}
    \label{fig:confidence_ellipse_1_1}
\end{figure}

\begin{figure}[htbp]
    \centering
    \includegraphics[width=1\textwidth]{confidence_ellipse_a=1,01.png}
    \caption{Доверительные эллипсы для $a=1.01$.}
    \label{fig:confidence_ellipse_1_01}
\end{figure}

\end{document}
